% Рисунок с доски: три базовых реологических элемента.
\begin{figure}[htbp]
\centering
\begin{tikzpicture}
  \pic at (0,0) {spring};
  \node[font=\footnotesize, anchor=north] at (0,-0.75)
    {упругий: $\Sdev = 2\mu\,\dev$};
  \node[font=\footnotesize, anchor=south] at (0,0.5) {закон Гука};

  \pic at (4,0) {damper};
  \node[font=\footnotesize, anchor=north] at (4,-0.75)
    {вязкий: $\Sdev = 2\eta\,\dt{\dev}$};
  \node[font=\footnotesize, anchor=south] at (4,0.5) {элемент Ньютона};

  \pic at (8,0) {stvenant};
  \node[font=\footnotesize, anchor=north] at (8,-0.75)
    {пластический: $N(\Sdev)\le\tau_s$};
  \node[font=\footnotesize, anchor=south] at (8,0.5) {элемент Сен-Венана};
\end{tikzpicture}
\caption{Базовые реологические элементы: пружина, демпфер и элемент сухого
трения. Из них методом реологических моделей собираются новые материалы.}
\label{fig:rheo-elements}
\end{figure}
